% Executive summary (round-2 revision, E1). Every number is a macro of the main text; the box restates results
% stated and qualified in the sections it points to.
\begin{center}
\fbox{\begin{minipage}{0.94\textwidth}\small\setstretch{1.0}
\textbf{Summary for operators and boards.}
\begin{enumerate}[leftmargin=*,itemsep=1pt,topsep=2pt]
\item \emph{The Rule of 40 is a pass mark, not a frontier.} Among the best software firms, a point of extra growth came with well under a point less margin: about a fifth of a point over 2010--25 on the GAAP basis, and \PsiLateOp{} in 2020--25 for growth of up to 40\%. These are choices across firms, not the trade-off any one firm faces, and the sign is not statistically firm (Section~\ref{sec:empirics}).
\item \emph{The basis decides the answer.} In 2023--25, \VBridgePassFcf\% of firm-years cleared 40 on free cash flow before stock-based compensation and \VBridgePassCash\% after it (Section~\ref{sec:data}).
\item \emph{Bought growth enters the score without its price.} For one serial acquirer with a 2025 GAAP score of \acqRGaapTwentyFive, the score on its cash margin before interest is \acqConvLo{} to \acqConvHi{} when its acquisitions are charged as capital, depending on how the charge is spread (Section~\ref{sec:acquired-growth}).
\item \emph{The top of the distribution fell.} The 95th percentile of the score fell by \TopFallOpEone{} points (GAAP) and \TopFallCashEone{} (free cash flow after stock-based compensation) between 2016--19 and 2023--25, while fewer firms grew fast; the GAAP decline survives a familywise correction at 10\% under most procedures and the cash decline under none. The data do not say why, and the top partly recovered in 2025.
\item \emph{Automation alone does not carry a firm past 100.} At the full-cost structure assumed here and fixed prices, complete automation keeps a mature firm's cash score below 100 unless its market grows faster than about \GstarCenOne\% a year (\GstarCcLo{} to \GstarCcHi\% over the credible range of compute cost), and \GstarPhiHalf{} or \GstarPhiSeven\% if half or 70\% of savings pass to customers. This is a bound at that cost structure, not a forecast or a target (Section~\ref{sec:horizons-ceiling}).
\item \emph{Execution sets the pace.} On a central scenario, a top-tail firm realizes about \PaceCenSix\% of the benefits of its automation program by 2036 and scores \CenRSix{} (\CenRSixLo{} to \CenRSixHi), between \CenRSixNoSM{} and \CenRSixAllSM{} depending on how much selling expense buys growth; execution, not market growth, moves the 2036 score most. Annual gains exceed the costs of transitions in flight after about \PayCrossCen{} years, but the whole program repays them only after \PayBackCen{} years (\PayBackDiscCen{} discounted), and its cumulative effect is negative at three and at five years (Section~\ref{sec:horizons-ranges}).
\item \emph{Which lever first depends on the spend basis.} Counting as growth spend only selling and marketing above a maintenance level, cutting general and administrative cost outranks automating the growth engine at \lpNGACutFirstNet{} of \lpNPositions{} observed positions; the growth engine leads only where selling spend is far above maintenance, and engineering ranks above both (Section~\ref{sec:levers-position}).
\item \emph{What lasts is what stays scarce.} Under free entry into a homogeneous product market, automation's savings lift the score only of a firm that holds an input rivals cannot replicate, and then only up to the full-cost ceiling (\autoref{prop:entry}).
\item \emph{Accountability is a requirement, not a headcount.} Internal control over financial reporting is pass-fail. In the paper's illustrative example no design built only from AI layers on one model provider passes, because their common failures are more than remote; whether that holds for a given firm depends on that probability, which is the number to estimate (\autoref{prop:amin}). The bottom-up accountability floor of a firm with \$1 billion of revenue is about \AminFloor{} to \AminFloorOne\% of revenue at the two thresholds (\AminFloorLo{} to \AminFloorHi\% over the range of audit fees).
\end{enumerate}
The paragraph ``How not to use these numbers'' at the end of Section~\ref{sec:intro} lists the misreadings to avoid; a glossary of terms is in the Supplementary Appendix (Section~\ref{supp:glossary}).
\end{minipage}}
\end{center}
